\chapter*{\centerline{Abstract}}
\graphicspath{{figs/Abstract/}}
\markboth{\MakeUppercase{Abstract}}{}
\iftoggle{fulltoc}{
  \addcontentsline{toc}{chapter}{Abstract}
}{}

Biomimetic robots offer a platform for studying neuromechanical control through experiments that would be impractical or unethical with living subjects. This dissertation advances a long-running effort at Portland State University to build a humanoid robot, from the trunk down, actuated by pneumatic artificial muscles and controlled by a synthetic neural network. The isometric force of Festo braided pneumatic actuators (BPAs) was characterized over a wide range of resting lengths, contractions, and pressures using a custom test jig. The analysis produced novel equations for the maximum isometric force as a function of resting length and a three-coefficient nondimensional force surface that predicts isometric force from relative strain and relative pressure. The actuator model was carried to the joint level through isometric knee torque experiments on two test stands: a simplified pinned knee and a biomimetic knee with a migrating instantaneous center of rotation. This approach isolated the sources of prediction error, revealing constant length offsets, series compliance in brackets and artificial tendons, and a loss of usable actuator length when BPAs wrap around the joint. The biomimetic knee driven by optimized BPA placements was shown to meet or exceed human isometric knee torque over most of the range of motion. Tools were created for curating relevant neuromechanical research, as well as simulating neuromechanical models based on human and robot biomechanical models.
